% ============================================================================
% DRAFT v2 -- Section 2: Related work
% Genre: benchmark / evaluation study. Subsection 2.3 places the paper in that
% genre explicitly. All \cite keys verified to exist in references.bib.
% ============================================================================

\section{Related Work}
\label{sec:related_work}

\subsection{Unsupervised Anomaly Detection}
\label{subsec:related_ad}

Unsupervised anomaly detection models normality and flags deviations, because
defect samples are scarce and the failures that matter are the ones not yet seen.
Four families dominate. Reconstruction methods learn to reproduce normal images
and score by reconstruction error, from autoencoders scored with a structural
similarity measure~\cite{bergmann2018improving} to generative
adversarial~\cite{akcay2019ganomaly,schlegl2019f} and memory-augmented
formulations~\cite{gong2019memorizing}. Feature-embedding methods characterise
the distribution of pre-trained deep features, either parametrically as in
PaDiM~\cite{defard2021padim}, by nearest-neighbour matching as in
SPADE~\cite{cohen2020spade} and PatchCore~\cite{roth2022patchcore}, or by density
estimation with normalizing flows~\cite{yu2021fastflow}. Teacher--student methods
distill a frozen teacher into a student that reproduces it on normal data, using
multi-scale feature matching~\cite{wang2021stpm,bergmann2020uninformed} or, in
reverse distillation, a decoder that reconstructs the teacher's features from a
bottleneck embedding~\cite{deng2022reverse}; the present paper studies a pipeline
of this last kind. Discriminatively trained methods inject synthetic
defects~\cite{zavrtanik2021draem} and have produced the strongest recent
results~\cite{liu2023simplenet,batzner2024efficientad,guo2025dinomaly,lee2022cfa};
surveys of the area are available~\cite{liu2024deep}.

Two properties of this literature matter here. The architecture is treated as a
given --- either the backbone of a pre-trained feature
extractor~\cite{defard2021padim,roth2022patchcore} or a fixed encoder--decoder
pair~\cite{deng2022reverse} --- and the step that turns a per-pixel anomaly map
into the single number on which a part is accepted or rejected is usually left
implicit. The pipeline we study reduces the multi-scale discrepancy map by an
unweighted mean over all pixels; Section~\ref{sec:leverage} shows that this one
undocumented choice is worth several times more accuracy than the architecture,
and that it can be revisited without retraining.

\subsection{Neural Architecture Search and the Evaluation of Its Search Phase}
\label{subsec:related_nas}

Neural architecture search automates the choice of network topology, from
early evolutionary~\cite{xie2017genetic} and reinforcement-learning
formulations~\cite{zoph2018learning,liu2018progressive} through differentiable
relaxation~\cite{liu2019darts}, platform-aware multi-objective
search~\cite{tan2019mnasnet} and regularised
evolution~\cite{real2019regularized}; see~\cite{ren2021comprehensive} for a survey.
Two lines of work are directly relevant.

The first applies NAS to anomaly detection. For image anomaly localisation,
Qian et al.~\cite{qian2024knowledgeguided} generate anomalous samples with a
knowledge-guided generator and use them as a reinforcement-learning reward, so
that architectures can be compared on a validation set that would otherwise be
unavailable. Work in this line takes for granted that a better architecture will
produce a better detector, and measures success by the downstream result alone.

The second line asks whether the search phase itself deserves that trust, and it
is the immediate predecessor of this paper. Yu et al.~\cite{yu2020evaluating}
compared the architectures returned by several NAS algorithms against
architectures sampled uniformly at random from the same search space, and found
that, on average, the NAS algorithms perform \emph{similarly to the random
policy}; they attribute this to weight sharing degrading the ranking of
candidates so that it no longer reflects true performance. Our findings agree in
direction and go further in three respects. First, from ``similar to random'' to
``worse than random'': in Section~\ref{subsec:search_vs_random} twelve independent
runs of the search all return one of the two configurations prone to a degenerate
evaluation, and the mean accuracy of the architectures they return sits at the
$1.2$th percentile of the distribution obtained by uniform sampling ($80.91\%$
against $81.65\%$, with the best backbone at $83.34\%$). Second, a mechanism
rather than a symptom: the cause is not weight sharing but an interaction between
the objective and the selection rule, as the evaluation can silently return
exactly $0.0$ when every validation batch is rejected, and since the search
minimises, that value is the optimum of the search's own range
(Section~\ref{subsec:proxy_stability}). Third, structure rather than signal
alone: two thirds of the search space cannot reach the deployed model at all
(Section~\ref{subsec:inert}), and the quantity being ranked contributes $0.54\%$
of the variance in accuracy against $95.32\%$ for the category
(Section~\ref{subsec:capacity_null}). Where Yu et al.\ study image classification
and language modelling with weight-sharing supernets, we study single-class
unsupervised anomaly detection with an independent quick-training proxy, and we
additionally show that the backbone ordering does not transfer between categories
($\bar\rho = 0.215$; $31\%$ of category pairs reverse), which makes a
single-category search unsound regardless of the size of the architecture effect.

\subsection{Benchmark and Evaluation Studies}
\label{subsec:related_benchmark}

The present paper belongs to a class of work whose contribution is a measurement
about how a field evaluates itself, rather than a new method. The class is
established, and we adopt three of its methods.

Bouthillier et al.~\cite{bouthillier2021accounting} showed that machine-learning
benchmarks routinely compare algorithms on a single run, that the unacknowledged
run-to-run variance is then large enough to reverse conclusions, and that the
number of runs a comparison requires can be derived from the variance. Their
analysis is the methodological template for
Section~\ref{subsec:resolvability}: we estimate the residual variance of the task,
fold in the externally reported effect size, and report the smallest difference
the standard design can resolve. IM-IAD~\cite{xie2024imiad} performed the
corresponding consolidation for industrial image anomaly detection, reproducing
published methods under one protocol and showing how much of the reported spread
survives standardisation. For anomaly detection outside the image domain the same
re-examination has been carried out on streaming
data~\cite{cao2024revisiting} and, most recently, on industrial multivariate time
series, where unified reproduction was combined with an event-level
evaluation~\cite{yu2026eventlevel}.

Our paper extends this line in two ways that we have not seen for unsupervised
image anomaly detection. It supplies the fully crossed configuration $\times$
category comparison that a variance decomposition requires --- 210 detectors,
every backbone of a search space on every category of the benchmark, one protocol
--- which is what turns a diagnosis of the architecture axis into a measurement
rather than an anecdote. And it anchors the resulting noise floor to the margin
the field actually contests, so that the resolvability of the standard benchmark
design can be stated in the same units as the claims it is used to support.
Neither IM-IAD nor the time-series reproductions reports that quantity; for this
task it turns out to be decisive.

\subsection{Knowledge Distillation}
\label{subsec:related_kd}

Knowledge distillation transfers the behaviour of a strong model to a smaller one
by matching outputs or intermediate representations~\cite{gou2021knowledge},
including similarity-preserving~\cite{tung2019similarity}, structured~\cite{liu2019structured}
and self-distilled~\cite{zhang2022self} variants and applications to
detection~\cite{wang2022distilling,yang2023attention}. In anomaly detection,
distillation is predominantly a \emph{detection} device: the teacher--student
discrepancy is the anomaly
signal~\cite{wang2021stpm,deng2022reverse,bergmann2020uninformed}. Using it as the
\emph{search metric} --- ranking candidates by how well a student distils from a
teacher, rather than by the accuracy of the detector it will eventually become ---
has been explored for NAS generally~\cite{kang2020towards} and, for anomaly
localisation, by Qian et al.~\cite{qian2024knowledgeguided}.

That reuse is the premise our results call into question.
Section~\ref{subsec:proxy_stability} evaluates the distillation proxy directly, by
repeating the search's own evaluation ten times per configuration, and finds that
it separates the twelve well-behaved backbones by only $5.7\%$ of its median ---
the same order as its own run-to-run noise --- while collapsing to an exact zero
for two of them. Its ordering is uncorrelated with the true ordering of the same
backbones: Spearman $\rho = -0.12$ over all fifteen categories and $\rho = 0.02$
on the very category the search is run on. The proxy is measuring something, but
not the thing it is used to select for.

\subsection{Positioning}
\label{subsec:related_summary}

This paper is a benchmark-and-evaluation study of a deployed design, not a new
detector, and we claim no priority for combining knowledge distillation with
architecture search. It differs from the NAS-for-anomaly-detection literature in
asking whether the search phase works rather than assuming it; it extends the
evaluation framework of Yu et al.~\cite{yu2020evaluating} to unsupervised anomaly
detection, quantifying the failure rather than merely detecting it; it adopts the
variance discipline of Bouthillier et al.~\cite{bouthillier2021accounting} and the
protocol consolidation of IM-IAD~\cite{xie2024imiad}, and reports the
resolvability of the benchmark that both leave unstated; and it contributes a
controlled comparison of a complete search space under one protocol
(Section~\ref{sec:positioning}) which makes the diagnosis measurable. On the
constructive side, it identifies the inference-time scoring policy as the axis
that actually moves accuracy and gives a label-free procedure for selecting it
(Section~\ref{subsec:pseudo_results}).
